% Teil II: durchgerechneter Weg, Teile a--d.

\subsection*{Gegeben und gesucht}
Massenpunkt im Punkt A (Hoehe $h_A$, Geschwindigkeit $v_A$) rollt
reibungslos; B tiefster Punkt (Hoehe $h_B$); C trifft eine Feder
(Federkonstante $k$); D höchster Punkt; Nulllage der $x$-Achse = rechter
Endpunkt der entspannten Feder. Gesucht: (a) $v_B$, (b) $v_C$, Federweg
$x$ und geleistete Arbeit, (c) Bedingung dafür, dass D erreicht wird,
(d) Skizze von $E_{\rm pot}(x)$.

\subsection*{Ansatz}
Energieerhaltung zwischen den Bahnpunkten; an der Feder wird kinetische
in Federenergie umgewandelt.

\subsection*{Rechnung}
\textbf{a) Geschwindigkeit in B.} Zwischen A und B (reibungsfrei):
\begin{equation}
  \tfrac12m v_A^2+mg\,h_A=\tfrac12m v_B^2+mg\,h_B
  \;\Rightarrow\;
  \boxed{\;v_B=\sqrt{v_A^2+2g(h_A-h_B)}\;}.
\end{equation}

\textbf{b) An der Feder (Punkt C).} Energie A $\to$ C (Hoehe $h_C$):
\begin{equation}
  v_C=\sqrt{v_A^2+2g(h_A-h_C)}\;.
\end{equation}
Beim Eindrücken der Feder wird diese Energie vollständig in
Federenergie umgewandelt (Ende der Fahrt kurzzeitig $v=0$):
\begin{equation}
  \tfrac12 k\,x^2=\tfrac12 m v_C^2
  \;\Rightarrow\;
  \boxed{\;x=v_C\sqrt{\frac{m}{k}}\;}.
\end{equation}
Die Arbeit des Wagens an der Feder ist gleich dieser Energie:
\begin{equation}
  W=\tfrac12 m v_C^2=\tfrac12 k x^2\;.
\end{equation}

\textbf{c) Ueberwinden des Punkts D.} Auf dem Rückweg aus der
Federenergie muss der Wagen die Hoehe $h_D$ erreichen; die Bedingung ist
(in Energie: am tiefsten C von der Feder zu D):
\begin{equation}
  \tfrac12 m v_C^2 \ge mg(h_D-h_C)
  \;\Longleftrightarrow\;
  v_C^2\ge 2g(h_D-h_C),
  \quad\text{d. h.}\quad
  x\ge\sqrt{\frac{2m(h_D-h_C)}{k}}\,.
\end{equation}
Gelingt gleicher Energiefluss bis D, so erreicht er D.

\textbf{d) $E_{\rm pot}(x)$.} Für $x<0$ (links der Feder) gibt es nur die
Gravitationsenergie, die dem Bahnprofil folgt (ansteigend zum
höchsten Punkt). Für $x\ge0$ (im Bereich der Feder) kommt die
Federenergie dazu,
$E_{\rm pot}=mgh(x)+\tfrac12 kx^2$, eine ansteigende
Parabel aus dem Kontaktpunkt. Die Gesamtkurve ist stetig und hat ihr
Minimum am tiefsten Bahnpunkt B.

\begin{ergebnis}
$v_B=\sqrt{v_A^2+2g(h_A-h_B)}$; $v_C=\sqrt{v_A^2+2g(h_A-h_C)}$,
$v$ bestimmt den Federweg $x=v_C\sqrt{m/k}$ und die Arbeit
$\tfrac12kx^2$; Bedingung $v_C^2\ge2g(h_D-h_C)$ (a. a. Federweg).
$E_{\rm pot}(x)$: Gravitation $+$ (nur rechts) Feder-Parabel.
\end{ergebnis}

\begin{kontrolle}
\textbf{Grenzfall:} Für $h_A=h_B$ (gleiche Hoehe) ist $v_B=v_A$ -- keine
Energieänderung, plausibel. \checkmark

\textbf{Einheiten:} $\sqrt{2g(h_A-h_B)}$: [m/s$^2$]$\cdot$m = m$^2$/s$^2$,
Wurzel [m/s]. $x=\sqrt{m/k}\cdot v$: $\sqrt{\mathrm{kg}\cdot\mathrm{m/N}}=\sqrt{\mathrm{m\cdot s^2/(s^2)}}\,$m? richtige Dimensionsprobe:
$[v\sqrt{m/k}]=(\mathrm{m/s})\sqrt{\mathrm{kg\,m/(kg\,m/s^2)}}=
(\mathrm{m/s})\sqrt{\mathrm{s^2}}=\mathrm m$. \checkmark

\textbf{D-Kriterium:} für $h_D=h_C+\tfrac{v_C^2}{2g}$ (genau erreichbar)
ist die Ungleichung gerade Gleichheit. \checkmark
\end{kontrolle}

\begin{merksatz}
Ohne Reibung bleibt $E=\tfrac12mv^2+mgh$ konstant; an der Feder gilt die
zusätzliche Umwandlung $\tfrac12mv^2\leftrightarrow\tfrac12kx^2$. Die
erreichbare Hoehe (auch nach Abprall) hängt nur von der zur Verfügung
stehenden kinetischen Energie, nicht vom Weg ab.
\end{merksatz}
